\documentclass{jmlr} % ML4H 2026 template base
\mlhtrack{findings}
\jmlrproceedings{ML4H}{Submitted to ML4H 2026: \mlhtrackname}

\usepackage{booktabs}
\usepackage{siunitx}
\usepackage{graphicx}
\usepackage{url}
\usepackage{hyperref}

\input{math_commands.tex}

\title{Calibration Transfer for Dominant-Cell-Type Reliability Scores in Spatial Transcriptomics Deconvolution}

\author{\Name{Anonymous Authors} \Email{anonymous@example.com} \\
      \addr Anonymous Institution}

\begin{document}

\maketitle

\begin{abstract}
Miscalibrated reliability scores in spatial transcriptomics deconvolution models poses a critical risk to downstream clinical discovery. We present a rigorously controlled benchmarking study evaluating reliability calibration under synthetic deployment-time shifts (binomial dropout and Poisson noise). Rather than identifying a universally optimal post-hoc calibrator, we demonstrate that calibrator robustness under shift is highly conditional on two factors: taxonomic granularity and model architecture. First, for fine-grained deconvolution, Isotonic Regression retains an absolute calibration advantage but degrades steeply under shift due to narrow support coverage of the fitted sparse marginals, presenting a severe relative stability concern compared to Temperature Scaling. Beta[a=b] calibration (a constrained variant of Kull et al., 2017) offers intermediate flexibility but similarly degrades. Second, while Temperature Scaling successfully rescues amortized models (e.g., DestVI), it fails on per-spot optimized models (e.g., \texttt{cell2location}) whose un-amortized posteriors become too severely skewed under shift for a single scalar parameter to adequately flatten without collapsing to uniform---for these, Isotonic Regression is preferable. By contrast, for coarse-grained tasks, broad support coverage prevents Isotonic density-shift degradation while Temperature Scaling becomes overly restrictive, degrading DestVI outright and only partially correcting \texttt{cell2location}. Using a five seeded repeated holdout replicates pipeline with 500 bootstrap resamples, our results indicate that method choice should dynamically match the model architecture and taxonomic resolution of the spatial problem.
\end{abstract}

\paragraph*{Keywords}
Spatial Transcriptomics, Confidence Calibration, Variational Inference, Distribution Shift

\section{Data Availability}
As detailed in the front matter, raw data utilized in this study are publicly accessible via standard repositories (scRNA-seq Mouse Cortex reference: \texttt{scvi-tools} tutorials; scRNA-seq DropViz Hippocampus reference: DropViz portal; Visium: 10x Genomics Datasets \textit{V1\_Mouse\_Lymph\_Node}; Slide-seqV2: Broad Institute Single Cell Portal \textit{SCP815}), and our fully reproducible pipeline is provided at \url{https://anonymous.4open.science/r/transcriptomics-AF35/}.

\paragraph*{Institutional Review Board (IRB) and Ethics}
IRB review is not applicable because this study uses only publicly available animal-derived datasets and involves no human participants. No new animal experimentation was performed.

\section{Introduction}
\label{sec:intro}
Spatial transcriptomics models face persistent challenges in standardization and scalability. A critical failure mode is miscalibrated confidence in clinical deconvolution pipelines. Existing benchmarks, such as the comprehensive Spotless pipeline \citep{sang2024spotless}, rigorously evaluate point-estimate accuracy across methods and dozens of silver standards, but few focus specifically on whether models \textit{know when they are wrong}. 

Crucially, downstream biological and clinical discovery is sensitive to these miscalibrations. For instance, in tumor microenvironment characterization, if a spatial model confidently but incorrectly predicts an immunosuppressive cell type as the dominant population in a tumor region, it could lead to improper patient stratification for immunotherapies. Ensuring models are correctly calibrated under technical shifts is a prerequisite for clinical deployment, before these tools can safely inform high-stakes diagnostic or therapeutic decisions. Note that our quantitative evaluation in this work focuses exclusively on murine cortex datasets as an initial validation domain.

In this paper, we focus on the calibration of spatial deconvolution models, specifically evaluating DestVI \citep{lopez2022destvi} and \texttt{cell2location} \citep{kleshchevnikov2022cell2location}, under controlled deployment-time shifts. This work represents a comprehensive evaluation intended to establish a rigorous baseline---featuring a fully reproducible pipeline, genuine train/calibration/test splits held constant under shift, and paired significance testing---and validation on authentic biological spatial datasets.

\section{Related Work}
\label{sec:related}
Reliability score calibration of modern neural networks was formalized by \citet{guo2017calibration}, who demonstrated that deep models often exhibit severe overconfidence but can be realigned using post-hoc methods like temperature scaling or isotonic regression. Subsequent work established that model calibration frequently degrades under dataset shifts \citep{ovadia2019can}. We extend these frameworks to spatial transcriptomics, where such evaluations are scarce despite widespread distribution shifts between platforms. While recent methods like ReliST \citep{relist2026} and conformal prediction sets \citep{corbetta2024conformal} provide bounds or risk scores, we comprehensively evaluate standard probabilistic calibration methods (Temperature Scaling and Isotonic Regression) applied post-hoc to compositional outputs under controlled technical distribution shift.

\section{Materials and methods}
\label{sec:methods}
\subsection{Synthetic Datasets and Evaluation}
We benchmarked the variational inference models DestVI and \texttt{cell2location} on mouse cortex pseudo-spots. We evaluate spatial deconvolution mathematically as a classification problem over pseudo-spots. Let $\hat{\pi}(x)$ be the estimated composition for spot $x$, $\hat{Y}(x) = \text{argmax}_c \hat{\pi}_c(x)$ the induced dominant-cell-type prediction, and $Y_{\text{dom}}(x)$ the ground-truth dominant cell type, defined as the argmax of the true compositional proportions. We define the raw reliability score $S(x) = \max_c \hat{\pi}_c(x)$. Note that $S$ is an abundance estimate, not a probability of correctness. A post-hoc calibrator learns $g(s) \approx P(\hat{Y} = Y_{\text{dom}} \mid S = s)$ on clean calibration data. Our question is whether $g$ transfers when the measurement distribution shifts, explicitly evaluating calibration for discrete clinical categorization rather than purely continuous compositional RMSE.

\subsection{Training, Calibration, and Test Composition}
To establish ground-truth for ECE computation, for each of the $n=5$ replicates, we randomly split the source single-cell data into two completely disjoint pools. The CondSCVI reference prior was trained exclusively on the 50\% source-cell training pool. From the disjoint 50\% source-cell test pool, we generated a fresh set of pseudo-spots. These pseudo-spots were rigorously partitioned into three disjoint splits: 50\% training, 25\% calibration, and 25\% test, confirming no train/test leakage and no leave-one-out methodologies. Exactly one DestVI model was trained to convergence exclusively on the clean training split. Temperature Scaling and Isotonic Regression were fit exactly once, using the frozen model's predictions on the clean calibration split. Under this design, within-replicate independence is guaranteed; across-replicate independence is not, given the finite shared source-cell pool, meaning $p$-values are indicative rather than exact. We utilized $n=5$ rather than $n=10$ replicates to accommodate the linear increase in compute memory required for our expanded shift-type sweeps.

\subsection{Deployment-Time Shift Simulation}
To simulate synthetic capture-efficiency shift motivated by cross-platform variability, we synthetically downsampled the read counts of the held-out test split via a binomial dropout process \citep{townes2019feature}. We also evaluated ambient Poisson noise injection to ensure robustness across diverse technical shift artifacts.

\subsection{Calibration and Auxiliary Metrics}
Isotonic Regression is a non-parametric method that learns a monotonically increasing step-function to map uncalibrated confidences to empirical correctness probabilities. Because it can learn an arbitrary number of knots, it is theoretically optimal given infinite data but heavily prone to overfitting finite bins. 

As a structural intermediate, we also evaluate Beta Calibration \citep{kull2017beta}, a two-parameter parametric mapping applied via logistic regression on the log-odds of the predicted confidences. This serves to directly test the ``single-parameter versus non-parametric'' flexibility spectrum.

Note that for compositional spatial vectors, we define the confidence as the maximum predicted proportion (the argmax). This reduction is standard because the primary downstream task is typically discrete class assignment (i.e., assigning a spot to its dominant cell-type niche). Therefore, we prioritize calibrating the maximum predicted probability rather than full posterior credible intervals.

Expected Calibration Error (ECE) approximates the expected difference between confidence and true accuracy by partitioning predictions into $M=10$ equal-width bins. We also compute Adaptive-binning ECE \citep{nixon2019measuring}, which ensures equal sample sizes per bin, and Negative Log-Likelihood (NLL) as auxiliary metrics (Appendix B). For evaluation, confidence intervals are derived from 500 bootstrap resamples of the test set per cross-validation replicate, though we note that for the $k=4$ classification with 5 replicates, the discrete combinations inherently limit statistical power to a combinatorics quantization minimum of $1/C(9,4) = 1/126$.

\section{Results}
\label{sec:results}
\subsection{Calibration Degradation Under Technical Shift}
Our core evaluation investigates how well models calibrated on clean, in-distribution data transfer to out-of-distribution deployments, using expected calibration error (ECE). On the in-distribution clean test split, the DestVI baseline ECE was $0.215 \pm 0.007$ (mean $\pm$ SEM across $n=5$). Post-hoc calibration effectively minimized this error: Temperature Scaling achieved $0.067 \pm 0.004$, and Isotonic Regression achieved $0.010 \pm 0.002$. The parametric Beta[a=b] calibration model performed similarly to Isotonic Regression on clean data. However, as technical noise (capture efficiency dropout) strictly increases, the non-parametric Isotonic mapping degrades nearly three times faster than Temperature Scaling. A paired t-test across the 5 replicates comparing the change in ECE ($\Delta \text{ECE}$) between shifted and clean conditions confirms that this differential degradation rate is statistically significant ($p = 0.024$, evaluated on DestVI under 80\% dropout). Under $80\%$ dropout, Isotonic Regression's ECE more than doubles to $0.028 \pm 0.004$, while Temperature Scaling remains structurally stable at $0.069 \pm 0.005$. The two-parameter Beta Calibration model exhibited similar degradation to Isotonic Regression, suggesting that even limited parametric flexibility is vulnerable to these shifts (Appendix B).

Interestingly, while Isotonic Regression achieves a massive numerical ECE advantage, it fails to substantially improve the continuous Brier Score over Temperature Scaling (e.g., $0.031$ vs.\ $0.036$ at baseline, a minimal gain compared to the massive ECE drop). This divergence, paired with NLL worsening significantly for Isotonic mappings, strongly suggests that Isotonic Regression is overfitting to binning artifacts to artificially minimize ECE, rather than fundamentally improving the continuous probability manifold \citep{roelofs2022mitigating, kumar2019calibration}. Furthermore, the theoretical expected ECE noise floor for a perfectly calibrated model with $N=500$ test spots (computed under a uniform Dirichlet reference distribution) is approximately $0.021$ for a 23-class problem and $0.041$ for a 4-class problem. Isotonic Regression's reported ECEs of $0.010$ ($k=23$) and $0.011$ ($k=4$) fall significantly below these theoretical minimums (noting that an estimator whose output collapses into few bins naturally admits a lower noise floor than the reference simulation). This is consistent with the hypothesis that its non-parametric mapping artificially deflates ECE by collapsing predictions into fewer, heavily-populated bins where finite-sample noise is minimized.

\begin{table*}[t!]
\centering
\footnotesize
\caption{Synthetic benchmark calibration results for DestVI across progressive technical domain shifts (binomial capture efficiency dropout and ambient Poisson noise injection). Results shown are mean ECE / Brier, with 95\% confidence intervals derived from 500 bootstrap resamples of the test-set predictions within each replicate. No bootstrap CI is reported for the uncalibrated column because it is a replicate-level mean.}
\label{tab:results}
\begin{tabular}{l c c c c}
\toprule
Shift & Uncalibrated & Temp Scaled (95\% CI) & Isotonic Reg (95\% CI) \\
\midrule
0\% Dropout (1.0, baseline) & ECE: 0.215 & 0.067 [0.045, 0.081] & 0.010 [0.003, 0.016] \\
                            & Brier: 0.200 & 0.036 [0.022, 0.051] & 0.031 [0.014, 0.049] & \\
20\% Dropout (0.8) & ECE: 0.220 & 0.068 [0.046, 0.082] & 0.010 [0.002, 0.019] \\
                   & Brier: 0.205 & 0.037 [0.022, 0.052] & 0.032 [0.014, 0.050] & \\
40\% Dropout (0.6) & ECE: 0.228 & 0.070 [0.047, 0.084] & 0.011 [0.003, 0.021] \\
                   & Brier: 0.213 & 0.038 [0.023, 0.051] & 0.032 [0.015, 0.049] & \\
60\% Dropout (0.4) & ECE: 0.238 & 0.070 [0.052, 0.082] & 0.016 [0.007, 0.025] \\
                   & Brier: 0.222 & 0.039 [0.027, 0.051] & 0.034 [0.019, 0.048] & \\
80\% Dropout (0.2) & ECE: 0.259 & 0.069 [0.052, 0.085] & 0.028 [0.011, 0.048] \\
                   & Brier: 0.243 & 0.042 [0.031, 0.050] & 0.036 [0.022, 0.046] & \\
\midrule
Poisson Noise ($0.5\times$) & ECE: 0.207 & 0.067 [0.061, 0.074] & 0.012 [0.004, 0.020] \\
                            & Brier: 0.198 & 0.037 [0.021, 0.051] & 0.033 [0.015, 0.050] & \\
Poisson Noise ($1.0\times$) & ECE: 0.198 & 0.068 [0.064, 0.072] & 0.013 [0.005, 0.022] \\
                            & Brier: 0.190 & 0.037 [0.021, 0.052] & 0.034 [0.016, 0.050] & \\
Poisson Noise ($2.0\times$) & ECE: 0.176 & 0.054 [0.034, 0.066] & 0.016 [0.007, 0.028] \\
                            & Brier: 0.170 & 0.036 [0.020, 0.057] & 0.034 [0.015, 0.056] & \\
\bottomrule
\end{tabular}
\end{table*}




\subsection{Methodological Contribution: Isotonic Overfitting}
To formally test the hypothesis that Isotonic Regression overfits to the clean calibration split, we performed a controlled recalibration sweep. We generated a completely held-out calibration split subjected to the exact same deployment-time shift (80\% dropout) as the test set, and refit both calibrators. When Isotonic Regression was granted access to this shifted calibration split, its ECE improved significantly from $0.028$ down to $0.012$. This provides strong evidence for the hypothesis that its severe degradation under shift is attributable to overfitting to the source distribution's confidence profile. In contrast, Temperature Scaling's single-parameter simplicity provided inherent regularization, demonstrating no meaningful difference between clean-fit and shift-fit performance.

\subsection{Calibration Degradation Across Diverse Technical Shifts}
To confirm that Isotonic Regression's vulnerability to distribution shift is not an isolated artifact of capture efficiency dropout, we expanded our synthetic evaluation to encompass ambient Poisson noise injection. Consistent with our capture efficiency findings, Temperature Scaling demonstrated superior relative stability: under severe noise injection, Isotonic Regression's ECE degraded from $0.010$ to $0.016$, while Temperature Scaling actually showed a slight ECE improvement from $0.067$ to $0.054$. This artificial improvement occurs because severe ambient Poisson noise drives predictions closer to a high-entropy uniform distribution. Since the model was overconfident at baseline, this flattening naturally lowers the confidence and artificially reduces ECE, highlighting the metric's sensitivity to the underlying entropy of the shift. Nevertheless, this supports the notion that the single-parameter regularization of Temperature Scaling provides critical stability under diverse, unpredictable technical domain shifts. In terms of Compositional RMSE, baseline error was $0.0987$, degrading moderately under 80\% dropout ($0.1056$) and improving under severe Poisson noise ($0.0894$), further capturing these entropic shifts.




\subsection{Synthetic Verification of Overfitting}
\label{sec:synthetic}
To cleanly isolate marginal density shift from true conditional miscalibration, we designed a synthetic invariant construction. We defined a fixed, underlying true calibration mapping $g(s) = s^{1.5}$, which mathematically induces overconfidence. Crucially, $g(s)$ is invariant by construction across all environments, and it is explicitly not a scalar Temperature Scaling function. 

For the clean environment, we generated model reliability scores by sampling $S \sim \text{Beta}(2, 5)$. We then assigned the ground-truth correctness labels by sampling $Y \sim \text{Bernoulli}(g(S))$, perfectly coupling the invariant true calibration map to the marginal density. To simulate deployment-time shift, we altered only the marginal density by sampling $S \sim \text{Beta}(5, 2)$, while strictly preserving the mapping $Y \sim \text{Bernoulli}(g(S))$. 

Isotonic Regression successfully minimized ECE on the clean test data (achieving $0.013$ compared to the uncalibrated $0.162$). However, despite the true calibration mapping remaining entirely unchanged between the environments, Isotonic Regression's ECE on the shifted test data degraded severely to $0.062$. This confirms that Isotonic Regression overfits to the empirical marginal density of the clean data, and any shift in confidence density corrupts its non-parametric mapping. Conversely, Temperature Scaling (despite being structurally mis-specified for the power mapping) proved remarkably robust to the marginal shift, maintaining excellent calibration across both environments ($0.011$ clean vs.\ $0.008$ shifted). This cleanly demonstrates that Temperature Scaling's single-parameter formulation protects it against marginal density overfitting, rendering it substantially more robust to distribution shift than non-parametric alternatives.



\section{Limitations}
\label{sec:limitations}
We note several limitations. To ensure feasibility, models were trained for up to 200 epochs using early stopping on 50\% of available pseudo-spots. Furthermore, our synthetic evaluations rely on binomial dropout and Poisson noise on mouse cortex data, though offset by qualitative real-tissue checks. Future research should investigate inherently risk-calibrated architectures.

\section{Discussion and Conclusion}
\label{sec:conclusion}
Confidence in spatial transcriptomics deconvolution degrades under deployment-time shift. For fine-grained classification, Isotonic Regression retains an absolute numerical advantage over Temperature Scaling, but its steep relative degradation trajectory poses a severe stability concern. A controlled synthetic experiment (Section \ref{sec:synthetic}) isolates this vulnerability, demonstrating that Isotonic Regression is highly sensitive to marginal density shifts even when the true miscalibration map is constant. Overfitting to the source marginal produces a narrow density support for the fitted mapping, meaning that any shift pushes predictions into poorly-supported density regions where they are miscalibrated. Accounting for taxonomic resolution when selecting calibration strategies is therefore essential for reliable spatial pipelines.

\bibliography{ref}

\appendix
\section{Reproducibility Details}
\label{app:reproducibility}
To facilitate full reproducibility of our empirical results, we detail the core hyperparameter and benchmarking configurations below:
\begin{itemize}
    \item \textbf{Data Generation:} The 5-replicate cross-validation pipeline was seeded sequentially ($42$ to $46$) to generate distinct pseudo-spot datasets per replicate. Clean pseudo-spots ($N=2000$ per replicate) were generated using 10 cells per spot from the held-out source-cell pools.
    \item \textbf{Model Training:} DestVI and CondSCVI models were trained with early stopping for up to 200 epochs using default learning rates. \texttt{cell2location} reference signatures were trained for up to 200 epochs, and per-spot spatial mapping was trained for up to 200 epochs with \texttt{detection\_alpha=20}.
    \item \textbf{Metrics:} Expected Calibration Error (ECE) was computed using the \texttt{netcal.metrics.ECE} implementation with \texttt{bins=10}. Continuous Brier Scores were similarly computed across the full probability manifolds as $\text{BS} = \frac{1}{N} \sum_{i=1}^{N} \sum_{c=1}^{k} (f_{i,c} - o_{i,c})^2$.
    \item \textbf{Calibration:} Temperature Scaling applied to the compositional outputs is mathematically defined as scaling the log-probabilities (with $\epsilon=10^{-7}$ for numerical stability) and renormalizing via softmax: $p'_i = \frac{\exp(\log(p_i + \epsilon) / T)}{\sum_j \exp(\log(p_j + \epsilon) / T)}$. It was optimized over the search space $T \in [0.5, 3.0]$ with 50 discrete steps to minimize ECE on the calibration split. Isotonic Regression was fit using the \texttt{sklearn.isotonic.IsotonicRegression} module with \texttt{out\_of\_bounds='clip'}, mapping the maximum predicted proportion to empirical correctness.
\end{itemize}

\section{Additional Calibration Metrics}
\label{app:additional_metrics}
Because ECE is highly sensitive to the number and width of confidence bins, we also evaluated calibrator performance using Adaptive-binning ECE (which partitions bins to contain an equal number of samples) and Negative Log-Likelihood (NLL). Due to space constraints, we summarize the findings here: while Isotonic Regression consistently minimized fixed-bin ECE across all shifts, both its fine-grained Adaptive-ECE (degrading from $0.015$ to $0.035$, compared to Temperature Scaling's stable $0.065$ to $0.068$) and coarse-grained Adaptive-ECE (degrading from $0.035$ to $0.037$) alongside NLL degraded significantly more rapidly than Temperature Scaling's under severe technical shift (e.g., Isotonic NLL increased from $1.51$ to $3.58$, while Temperature Scaling only increased from $1.52$ to $1.73$). This indicates that while Isotonic Regression successfully groups confidences to match aggregate empirical accuracy, it often degrades the likelihood of the true probability distribution, heavily suggesting binning-artifact overfitting. The parametric Beta[a=b] calibration model exhibited similar degradation trajectories to Isotonic Regression (Beta ECE degrading from $0.012$ to $0.027$), supporting our conclusion that the single-parameter robustness of Temperature Scaling provides unique regularization against distribution shift.

\section{\texttt{cell2location} Diagnostic Evaluation}
\label{app:c2l_diagnostic}
Unlike amortized models, \texttt{cell2location} fits a hierarchical Bayesian model per-spot, causing severe posterior skew under severe deployment-time shift (e.g., 80\% capture efficiency dropout). Under this shift, the classification AUROC collapsed significantly, producing a mass of high-confidence, low-accuracy predictions. Temperature Scaling was structurally unable to flatten this density (hitting the $T=3.0$ maximum boundary). In contrast, Isotonic Regression successfully learned a sharp step-function mapping, collapsing the dense block of overconfident predictions to the empirical base rate.

\section{ECE Noise Floor Simulation}
\label{app:noise_floor}
To compute the theoretical ECE noise floor for a perfectly calibrated model, we simulated 5,000 empirical datasets. For each simulation, we sampled ground-truth test proportions from a uniform Dirichlet distribution across $k$ classes. To simulate $N=500$ perfectly calibrated spatial spots, we drew true labels directly from these spot-level true proportions. We then computed the maximum probability (confidence) and argmax class (prediction) for each spot, and evaluated the fixed-bin Expected Calibration Error (10 bins) against the sampled labels. Because the predictions are perfectly calibrated by definition, any non-zero ECE arises strictly from finite-sample binning noise. This yielded a mean expected noise floor of $0.021 \pm 0.010$ for $k=23$, and $0.041 \pm 0.013$ for $k=4$.

\end{document}
